\scanpage{30}
s'insérant dans le diagramme de suites exactes (auquel il donne naissance)

\diagram{\begin{tikzcd}[ampersand replacement=\&,column sep=0.35em,row sep=0.7em,cramped]
\&\&\&1\arrow[d]\&\&\&\\
\&1\arrow[dr]\&\&\mathfrak S_I^+\times\mathbb F_2^{\underline\omega}\arrow[dl,"\mathrm{pr}_1"']\arrow[dr,"\mathrm{pr}_2"]\arrow[dd]\&\&1\arrow[dl]\&\\
1\arrow[dr]\&\&\mathfrak S_I^+\arrow[dl,Rightarrow]\arrow[dr]\&\&\mathbb F_2^{\underline\omega}\arrow[dl]\arrow[dr,Rightarrow]\&\&1\arrow[dl]\\
{(23)}\&\mathfrak S_I^+\arrow[dr]\&\&\mathrm{Gl}_{\Lambda}(M)\arrow[dl]\arrow[dr]\arrow[dd,"{\text{cart}}" description]\&\&\mathbb F_2^{\underline\omega}\simeq V_{S_Q}\arrow[dl]\&\\
\&\&\mathfrak S_I\arrow[dl]\arrow[dr,"\mathrm{sg}"]\&\&\mathcal D_Q\arrow[dl,"\chi_C"']\arrow[dr]\&\&\\
\&1\&\&\{\pm1\}\arrow[dl]\arrow[dr]\arrow[dd]\&\&1\&\\
\&\&1\&\&1\&\&\\
\&\&\&1\&\&\&
\end{tikzcd}}

\uline{Proposition} Le \inserted{noyau \deleted{\unsure{de}}($\simeq\mathfrak S_I^+$) de l'}hom. canonique $\mathrm{Gl}_{\Lambda}(M)\to\mathcal D_Q$ \deleted{induit un iso} est aussi égal au groupe des commutateurs du sous-groupe \inserted{distingué} $\mathrm{Sl}_{\Lambda}(M)$ d'indice $2$ de $\mathrm{Gl}_{\Lambda}(M)$ (NB $\Lambda^*=(\mathbb Z/4\mathbb Z)^*=\{\pm1\}$), qui est \inserted{donc} d'indice $4$ dans $\mathrm{Sl}_{\Lambda}(M)$ -- le quotient $(\mathrm{Sl}_{\Lambda}(M))_{\mathrm{ab}}$ est \deleted{isomorphe à $\mathbb Z/4\mathbb Z$} \inserted{cyclique d'ordre $4$}, et l'hom

\[
\mathrm{Gl}_{\Lambda}(M)\xrightarrow{\det}\{\pm1\}
\]

est le composé \deleted{\unsure{des deux hom|des hom}}

\[
\mathrm{Gl}_{\Lambda}(M)\longrightarrow\mathcal D_Q
\underbrace{\xrightarrow{\deleted{\unsure{\chi}}\,\det}}_{\substack{\text{noyau }\mathcal D_Q^+,\\\text{cyclique d'ordre }4}}\{\pm1\}. \tag*{(24)}
\]

\uline{Corollaire 1} On a un sous-diagramme de suites exactes de (23), obtenu ainsi :
% archive sheet number: 29
